\documentclass[12pt]{article}
\usepackage{xeCJK}

\usepackage{charter}
\usepackage{fullpage}
\usepackage[colorlinks=false]{hyperref}
\usepackage{ifthen}
\usepackage{comment}
\usepackage[title,titletoc]{appendix}
\usepackage{pagecolor}
\usepackage{amsmath}
\usepackage{amsfonts}
%\usepackage[normalem]{ulem}
\usepackage{siunitx}
\usepackage{amsthm}
\sisetup{per=slash, load=abbr}

\usepackage{pgfplots}
\usetikzlibrary{positioning}
\usetikzlibrary{fit}
\usetikzlibrary{snakes}
\usetikzlibrary{shapes.geometric}
\usetikzlibrary{patterns}
\usetikzlibrary{shapes,arrows,chains}
\usepgfplotslibrary{patchplots,colormaps}
\usetikzlibrary{calc}
\usetikzlibrary{positioning, fit}
\usetikzlibrary{backgrounds}
\usetikzlibrary{intersections}

\newcommand{\whitepaper}[1]{\begin{center}\fbox{\parbox{0.75\textwidth}{{\small
#1}}}\end{center}}

\newcommand{\pcolor}{orange!25}

\usepackage{setspace}
\usepackage{algorithm2e}
\bibliographystyle{ieeetr}

\usepackage{geometry}
\geometry{left=3cm,right=3cm,top=1.6cm,bottom=3cm,headheight=0pt,headsep=1.5em}
\usepackage{fancyhdr}
\pagestyle{fancy}
\chead{\includegraphics[scale=0.2]{../common/Nebulas.png}}  %在此处插入logo.pdf图片 图片靠左
\lhead{} % 页眉中间位置内容
\rhead{}
%\setlength{\topskip}{1em}

\usepackage{indentfirst}


\setCJKmainfont[BoldFont = STSongti-SC-Bold]{STSongti-SC-Regular}
\setCJKfamilyfont{hei}{SIL-Hei-Med-Jian}		%宋体
\setCJKfamilyfont{song}{SimSun}		%宋体
\setCJKfamilyfont{kai}{Kaiti}		%楷体
\setCJKfamilyfont{fang}{song}	%仿宋
\setCJKfamilyfont{li}{song}			%隶书
\setCJKfamilyfont{you}{Yuanti}		%幼圆

\newcommand{\song}{\CJKfamily{song}}	%宋体
\newcommand{\hei}{\CJKfamily{hei}}	%黑体
\newcommand{\kai}{\CJKfamily{kai}}	%楷体
\newcommand{\fs}{\CJKfamily{fang}}	%仿宋
\newcommand{\li}{\CJKfamily{li}}		%隶书
\newcommand{\you}{\CJKfamily{you}}	%幼圆
\newcommand{\reffig}[1]{图\ref{#1}}
\newcommand{\refsec}[1]{\S \ref{#1}}

\onehalfspacing   % ----------设置1.5倍行距（可能有意义，待调整）

%\parindent=20pt  % -------------------首行缩进大小，英文分段就直接0pt了吧。
\setlength{\parindent}{2.1em}
\setlength{\parskip}{0.3\baselineskip}
\newcommand{\nrcore}{Core Nebulas Rank}
\newcommand{\nrext}{Extended Nebulas Ranks}
\newcommand{\dom}{{\; \texttt{dom}\;}}

\setCJKsansfont[BoldFont = STHeitiSC-Medium]{STHeitiSC-Light}


\newtheorem{property}{特征}
%\addbibresource{reference.bib}

\begin{document}
\pagestyle{empty}
\renewcommand{\contentsname}{目录}
\renewcommand{\abstractname}{摘要}
\renewcommand{\refname}{参考文献}
%\renewcommand{\nomname}{术语表（按首字母排序）}
\renewcommand{\figurename}{图}
\renewcommand{\tablename}{表}
\renewcommand{\baselinestretch}{1.5}
\renewcommand{\appendixname}{附录}
\renewcommand{\proofname}{证明}

\pagecolor{\pcolor}

\begin{titlepage}
  \begin{center}
    \vspace*{5.5cm}
    \includegraphics[scale=0.5]{../common/Nebulas.png}
    \vspace{0.5cm}


    \textbf{\huge{星云治理橙皮书}}

    \vspace{0.5cm}
    李晨、王冠
    \vfill
    2019年4月 \\
    版本号:0.0.3
    \textbf{}
  \end{center}

\end{titlepage}
\setcounter{page}{0}
%\thispagestyle{empty}
\tableofcontents
\newpage
\setcounter{page}{1}
\pagestyle{fancy}
\vspace*{0.01cm}
% !TEX root = main.tex
\section{概述}

\textbf{关键字： 治理\ 去中心化\ 协作\ 激励\ 自进化\ 自治 }

\vspace{2em}

星云链（Nebulas）是开源公链，星云是自治元网络（Autonomous Metanet）~\cite{AutonomousMetanet}，星云专注于处理复杂数据和交互、复杂的协作关系，致力于通过区块链等技术手段，实现让每个人从去中心化协作中公平获益的愿景~\cite{vision}。

本橙皮书着重介绍为了实现愿景，星云如何借助于独有的创新技术\textbf{治理链上公共资产，探索全新的去中心化协作模式，实现提供正向激励、可自进化的去中心化自治组织}（Decentralized Autonomous Organization，DAO）~\cite{DAO}。星云治理在如下三方面进行了独创设计，本书将一一介绍：

\begin{enumerate}
	\item \textbf{组织结构和监督机制：}
	独立运行且互相制约的星云三会：星云理事会、星云基金会、星云技术委员会，阐述其基本构成、权力和义务；
	\item \textbf{链上协作方式：}
	星云社区项目制介绍，星云通过“NAT链上投票”（参见\ref{nat}）实现社区协作和系统升级的流程；
	\item \textbf{经济体和激励：}
	在星云治理过程中，星云经济体如何对每一位社区成员提供正向激励。
\end{enumerate}
%\printbibliography
\newpage
\bibliography{reference}

\newpage
\begin{appendices}
\section{Change Log}
\begin{itemize}
\item{0.0.9} Release.
\item{0.1.0} Update Figures and polish.
\end{itemize}

\end{appendices}

\end{document}
